\section{Classical Cut Elimination with Explicit Multiplicities}
\label{sec:p07-classical-cut-elimination}

Cut elimination removes an intermediate formula from a derivation while
preserving its endsequent. The theorem here concerns exactly
\hyperref[sheet:lk-defneg]{\(\mathsf{LK}_{\mathrm d}\)}: classical
two-sided finite multisets, explicit weakening and contraction, and
negation defined by \(\neg A\defeq A\to\bot\).
The proof retains the distinction between a formula and a selected
occurrence of that formula. This distinction is essential when contraction
duplicates the material to be eliminated.

For comparison, Gallier gives a semantic cut-elimination route for his
\(LK\) and a constructive route for a separate set-context calculus in
negation normal form \cite[Secs.~6.1--6.2, 6.4.4]{Gallier2015}.
The following proof is a direct construction for the registered multiset
calculus; its contraction and implication cases are included explicitly.
Zach's generalized-calculus treatment gives a different all-occurrences
Mix formulation and a sketched quantifier extension
\cite[Thm.~10; Sec.~13]{Zach2021GeneralizedV2}.
Here the admissibility proof follows the selected occurrence tokens,
including any retained copies of the same formula, and supplies the
quantifier and exact-syntax realization arguments in full.

\subsection{Exact syntax and occurrence bookkeeping}

\begin{hbdefinition}{def:p07-classical-cut-domain}{The derivations being transformed}
Formulas use the primitives
\(\bot,\land,\lor,\to,\forall,\exists\), arbitrary finite-arity
predicate and function symbols, and the fixed infinite variable alphabet
of \cref{def:p02-fol-signature}. Equality is absent. Expand defined
negation before counting formula constructors.
The formula rank \(\operatorname{rk}_{\mathrm{fm}}(A)\) counts logical
symbols: a predicate atom has rank zero, \(\bot\) has rank one, a binary
connective adds one to the sum of its operands' ranks, and a quantifier
adds one to its matrix's rank. The tree height \(\operatorname{ht}(D)\)
is zero at an initial sequent and one plus the maximum premise height
at every other inference.

A sequent \(\Gamma\Rightarrow\Delta\) has finite multiset contexts.
Identity is exactly \(A\Rightarrow A\), without an implicit context.
The other initial schema is the contextual
\(\bot,\Gamma\Rightarrow\Delta\). Exchange is built into multiset
notation; weakening and contraction remain actual inferences.
The rule \(\to L\) splits its two pairs of contexts, whereas
\(\land R\) and \(\lor L\) share theirs.
All quantifier instances use the total capture-avoiding substitution in
\cref{def:p02-fol-capture-avoiding-substitution}. The eigenvariable of
\(\forall R\) or \(\exists L\) is not free in its conclusion.

A derivation is \emph{cut-free here} when it contains neither Cut nor Mix.
Only the registered initial sequents are permitted; additional theory
initial sequents, equality, induction or other logical rules are outside
this theorem. Arbitrary finite formula assumptions may still occur in
the antecedent contexts already allowed by the sheet.
\end{hbdefinition}

\begin{hbdefinition}{def:p07-classical-marked-mix}{Selected copies and the strengthened Mix task}
Write the two endsequents as
\[
 \Gamma\Rightarrow\Delta,A^m,
 \qquad A^n,\Pi\Rightarrow\Sigma,
\]
where \(A^m\) and \(A^n\) mark exactly the selected occurrences.
Every context may contain further, unselected copies of \(A\).
The task is to construct a cut-free derivation of
\(\Gamma,\Pi\Rightarrow\Delta,\Sigma\).
For the induction we allow \(m,n\geq0\); this is a strengthened
admissibility statement, not an extra zero-multiplicity rule on the sheet.
When both are positive it is exactly the registered Mix task.

Marks are auxiliary annotations on proof occurrences. Pull a context mark
back to its corresponding premise occurrence; in a shared context pull it
back to both premises. If contraction's retained occurrence is marked,
mark both contracted ancestors. If it is unmarked, leave both ancestors
unmarked. No formula occurrence is deleted just because it has the same
spelling as a selected occurrence.
\end{hbdefinition}

\begin{hbdefinition}{def:p07-classical-alpha-construction}{An auxiliary presentation for bound-variable names}
Use finite multisets of alpha-equivalence classes for an auxiliary proof
construction, retaining every rule, context multiplicity and side condition
of the registered sheet. The class relation is the syntactic
alpha-equivalence of \cref{def:p02-fol-alpha-equivalence}.
Here \(A=_{\alpha}B\) abbreviates the source's \(A\equiv_\alpha B\);
this notation does not assert elementary or semantic equivalence.
Write \([A]\) for the class of \(A\).
Quantifier instances are the classes of the specified total
capture-avoiding instances. Class identity has one copy on each side.

This auxiliary presentation makes bound-variable choices harmless during
the induction. Its cut-free output will be realized on the exact registered
sheet with prescribed literal endsequent representatives. No unproved
alpha-conversion inference is being added to that sheet.
\end{hbdefinition}

\begin{hblemma}{thm:p07-classical-alpha-identity}{proved=annex:p07.classical.alpha-identity}{Cut-free identity between alpha-equivalent representatives}
If \(A=_{\alpha}B\), the exact registered sheet has a cut-free derivation
of \(A\Rightarrow B\). The construction treats every connective and
both quantifiers and uses no semantic completeness theorem.
\end{hblemma}

\begin{hblemma}{thm:p07-classical-parameter-substitution}{proved=annex:p07.classical.parameter-substitution}{Parameters and term substitution in the auxiliary construction}
In the auxiliary presentation, an eigenvariable can be replaced by a fresh
parameter throughout its premise derivation while the conclusion of that
eigenvariable inference remains unchanged. Fresh parameters can avoid any
specified finite set.
Applying one capture-avoiding term substitution to a cut-free derivation's
endsequent has a cut-free auxiliary derivation of no greater height, after
freshening its eigenvariables as needed.

Formula rank, defined as the number of logical symbols after expanding
negation, is unchanged by term substitution or alpha-conversion.
\end{hblemma}

\begin{hblemma}{thm:p07-classical-alpha-realization}{proved=annex:p07.classical.alpha-realization}{Realizing the exact endsequent}
Every cut-free auxiliary derivation can be converted into an exact
\(\mathsf{LK}_{\mathrm d}\)-derivation with any prescribed representatives
of its endsequent classes, preserving all multiplicities. The resulting
derivation is cut-free; no preservation of its physical height is asserted.
Conversely, passage from an exact proof to classes preserves its inference
tree and height.
\end{hblemma}

The height claim belongs to the auxiliary construction. For example,
replacing just the right occurrence in the exact identity
\(\forall xP(x)\Rightarrow\forall xP(x)\) by
\(\forall yP(y)\) does not leave a literal identity leaf. The realization
lemma supplies its cut-free proof instead of incorrectly calling it the
same zero-height proof.

\subsection{The strengthened induction}

\begin{hblemma}{thm:p07-classical-bottom-erasure}{proved=annex:p07.classical.bottom-erasure}{Erasing selected right-hand falsities}
From a cut-free derivation of \(\Gamma\Rightarrow\Delta,\bot^m\), with
\(m\geq0\) selected copies, one can construct a cut-free derivation of
\(\Gamma\Rightarrow\Delta\). The result holds in both the auxiliary
presentation and the exact sheet. Unselected copies in \(\Delta\) remain.
\end{hblemma}

\begin{hbtheorem}{thm:p07-classical-marked-mix}{proved=annex:p07.classical.marked-mix}{Admissibility of the full marked Mix task}
Given cut-free derivations of the two sequents in
\cref{def:p07-classical-marked-mix}, for arbitrary \(m,n\geq0\), one
can construct a cut-free derivation of
\(\Gamma,\Pi\Rightarrow\Delta,\Sigma\) on the exact registered sheet.
The result eliminates only the selected copies, including when contraction
has supplied several ancestors for one selected end occurrence.
\end{hbtheorem}

\begin{hbmethod}{def:p07-classical-cut-construction}{A terminating local construction}
First pass to the auxiliary presentation. For cut-free input proof trees
\(D,E\), induct lexicographically on
\[
 \bigl(\operatorname{rk}_{\mathrm{fm}}(A),
       \operatorname{ht}(D)+\operatorname{ht}(E)\bigr).
\]
Use the ranks and heights in \cref{def:p07-classical-cut-domain}; here the
heights belong to the auxiliary trees before exact-syntax realization.

Zero multiplicity uses weakening. Selected structural rules reduce to their
premises, with both contraction ancestors selected. A nonselected logical
inference is commuted by applying the induction to each immediate premise
and restoring that inference; duplicated split contexts are contracted.
Before a quantifier permutation, freshen its eigenvariable away from the
other proof.

When both final logical inferences introduce selected copies, first remove
the residual selected copies of \(A\) from their immediate premises by
same-rank calls against the opposite \emph{original} proof. Those calls
have smaller original height sums. Then eliminate the exposed immediate
subformulas or quantified instances using smaller-rank calls. Their proof
heights may be larger: rank, rather than height, decreases at this stage.
Finally realize the exact endsequent representatives.
\end{hbmethod}

\begin{hbexample}{def:p07-classical-contraction-case}{Why a unit Cut is not the induction hypothesis}
If the first proof ends with right contraction of \(A,A\) into one selected
\(A\), the recursive task has two selected copies in that premise.
Trying to use only a unit-cut statement leaves one ancestor untreated.
The marked Mix lemma accepts the increased multiplicity while the premise
height decreases. Contraction of an unselected copy instead keeps both
ancestors unselected and is reapplied afterward.
\end{hbexample}

\begin{hbexample}{def:p07-classical-principal-conjunction}{A principal conjunction example}
Identity and weakening derive \(P,Q\Rightarrow P\) and
\(P,Q\Rightarrow Q\); \(\land R\) gives
\(P,Q\Rightarrow P\land Q\).
Identity \(P\Rightarrow P\), left-weakened by \(Q\), followed by
\(\land L\) gives \(P\land Q\Rightarrow P\).
A Cut between them concludes \(P,Q\Rightarrow P\).
Its cut-free replacement is already identity on \(P\) followed by left
weakening. The general principal-conjunction reduction in the annex also
handles residual selected copies and arbitrary contexts.
\end{hbexample}

\begin{hbexample}{def:p07-classical-eigenvariable-case}{Freshness is a rule condition}
From \(P(y)\Rightarrow P(y)\), one may not infer
\(P(y)\Rightarrow\forall xP(x)\) by \(\forall R\): \(y\) is free in
the proposed conclusion's antecedent. A two-element structure with \(P\)
true of the value assigned to \(y\) and false of the other element also
shows why that inference is unsound. Renaming an eigenvariable must rename
its premise derivation coherently; it does not waive the side condition.
\end{hbexample}

\subsection{Elimination and the cut-free consequences}

\begin{hbtheorem}{thm:p07-classical-cut-elimination}{proved=annex:p07.classical.cut-elimination}{Cut elimination for the exact defined-negation classical calculus}
Every finite derivation on \texttt{sheet:lk-defneg}, including registered
Cut and Mix instances, can be transformed into a derivation of the same
literal endsequent containing neither rule. Context multiplicities and
the defined-negation language are preserved.
\end{hbtheorem}

\begin{hbcorollary}{thm:p07-classical-analytic}{proved=annex:p07.classical.analytic}{Analytic derivations and the empty sequent}
Every derivable propositional sequent on this sheet has a cut-free
derivation in which each formula occurrence is a literal subformula of an
endsequent formula, after expanding defined negation.
For first-order cut-free proofs, every formula occurrence lies in the
closure of the end formulas under taking immediate subformulas,
capture-avoiding term instances and alpha-renaming. Literal subformulas
alone do not suffice.
The empty sequent has no derivation on this sheet.
\end{hbcorollary}

The construction proves no bound on the size of the resulting exact proof
and no finite first-order search bound. Midsequent permutations, Herbrand
extraction and calculus translations require their additional constructions.
The separate intuitionistic, primitive-negation and equality rule inventories
are not included in the theorem merely because they also use the name
``sequent calculus.''
